% !TeX root = ../../Book3.tex
\chapter{调和函数与次调和函数}
\label{b3:ch:15}
\BookChapterNumber{b3:ch:15}
\begin{chapterabstract}
全纯函数的实部满足 Laplace 方程，圆周平均则把这个微分方程转化为比较和估计。本章由平均值性质建立最大值原理、Poisson 表示、Harnack 控制与局部紧致性，再用周期说明调和共轭何时能够拼成单值函数。为了保留全纯函数零点处的对数奇性，随后把平均等式放宽为次调和不等式，并以正分布和 Riesz 测度记录非调和部分。Weyl 引理与局部对数势分解将微分方程、点态代表和测度联系起来，为后续边值问题提供工具。
\end{chapterabstract}

\begin{learninggoals}
\begin{itemize}
\item 从平均值性质恢复调和方程与光滑性，并用 Poisson 核推出内部估计和正调和函数的边界测度表示。
\item 检查 Harnack 比较中的正性和区域条件，说明局部有界性怎样产生紧致性，以及单调性怎样控制调和函数列。
\item 计算共轭微分形式的周期，区分单值调和共轭与全局全纯对数模的存在条件。
\item 在允许负无穷值的情形使用次调和比较、最大值原理及稳定运算，辨明上半连续性与平均不等式各自的作用。
\item 用测试函数计算分布 Laplace 算子，通过正性识别 Riesz 测度，并准确使用分布判据、Weyl 引理和局部 Riesz 分解。
\end{itemize}
\end{learninggoals}

函数 \(\log|z|\) 给出了进入本章的一条线索。在不含原点的小圆盘上，它是某个全纯对数的实部；绕原点一周，共轭函数却增加 \(2\pi\)，因而不能在整个去心平面上取成单值函数。若把原点处的值设为 \(-\infty\)，它仍局部可积，并满足次调和平均不等式；它的分布 Laplace 算子又在原点留下一个点质量。局部方程、全局周期和奇性的测度描述，需要分别建立。

\ProseParagraphBreak
\secref{b3:sec:1501}从圆周平均出发。Poisson 核既构造连续边值的调和延拓，也给出 Harnack 不等式、导数估计和局部紧致性。正调和函数的边界测度表示进一步容纳不连续甚至集中的边界数据；它也提示，几乎处处的径向边值并不总能确定一个调和函数。\secref{b3:sec:1502}再把调和函数接回全纯函数：方程保证局部共轭存在，闭曲线周期决定这些局部共轭能否拼合，而周期属于 \(2\pi\mathbb Z\) 已足以使指数单值。

\ProseParagraphBreak
\secref{b3:sec:1503}用上半连续性和平均不等式处理零点处的对数奇性。比较原理使这一概念可以在小圆盘上验证，也使最大值、拼接和极限运算有统一的依据。\secref{b3:sec:1504}转到分布语言：先确定 \(\Delta\log|z|\) 的常数，再由正性把 Laplace 分布识别为测度。分布只看见几乎处处等价类，因此恢复次调和函数时还须选定点态代表；局部 Riesz 分解中的处处等式也依赖这一步。

\ProseParagraphBreak
本章使用平面微积分、前面的全纯函数局部理论，以及第一卷的局部可积函数、近似单位和正泛函的测度表示。阅读时可对照 Axler、Bourdon、Ramey 的《Harmonic Function Theory》\cite{Axler2001Harmonic}第4--8页的均值与最大值原理、第12--18页的球上 Poisson 表示与均值逆命题，以及第47--50页的 Harnack 理论；比较维数变化时，应重新核对核函数及估计的常数。Ransford 的《Potential Theory in the Complex Plane》\cite{Ransford1995Potential}可作为平面势论的进一步阅读方向。

% !TeX root = ../../Book3.tex
\section{调和函数}
\label{b3:sec:1501}

全纯函数的实部不只是满足一条微分方程。它在每个圆盘中心都等于边界的平均值；反过来，这种平均性质足以恢复光滑性和方程。先把这两个描述联系起来，再从正性得到定量控制。

\subsection{Laplace 方程}
\label{b3:subsec:150101}

设 \(D\subseteq\mathbb C=\mathbb R^2\) 为开集。实值函数 \(u\in C^2(D)\) 若满足 \(\Delta u=u_{xx}+u_{yy}=0\)，称为\term[diaohe-hanshu]{调和函数}[harmonic function]。复值函数调和是指实部、虚部均调和；最大值和正性结论以下只用于实值函数。\(x\)、\(x^2-y^2\)、\(\log|z|\) 分别在全平面、全平面和去心平面上调和。最后一个例子提醒我们：定义域的缺口可能承载奇性，不能在原点直接补值。

由 Cauchy--Riemann 方程和全纯函数的光滑性，\(f=u+\mathrm{i}v\) 全纯时 \(u_{xx}+u_{yy}=v_{yx}-v_{xy}=0\)。实变量的方程在二维保留了复变量的局部结构，但尚未保证全局调和共轭存在。

\subsection{平均值性质}
\label{b3:subsec:150102}

\term[pingjunzhi-xingzhi]{平均值性质}[mean-value property]把方程变成与尺度有关的等式。
\begin{theorem}[调和函数的平均值性质]\label{b3:thm:15-mean-value}
若 \(u\) 在 \(D\) 调和，且 \(\overline{B(a,r)}\subset D\)，则
\[
u(a)=\frac1{2\pi}\int_0^{2\pi}u(a+r e^{\mathrm{i}t})\,\mathrm{d}t
=\frac1{\pi r^2}\int_{B(a,r)}u\,\mathrm{d}A.
\]
反之，连续函数在每个内部圆盘上满足圆周平均等式，则为调和函数。
\end{theorem}
\begin{proof}
圆周平均记作 \(m(r)\)。微分和散度公式给出
\[
 m'(r)=\frac1{2\pi r}\int_{\partial B(a,r)}\partial_nu\,\mathrm{d}s
 =\frac1{2\pi r}\int_{B(a,r)}\Delta u\,\mathrm{d}A=0.
\]
由连续性 \(m(r)\to u(a)\)，得第一式；乘 \(2r\) 并积分得面积版本。
反向证明不能预先假定函数可微。在相对紧的小开集上取光滑径向非负核 \(\eta_\varepsilon\)，积分为一。极坐标和假设给 \(u*\eta_\varepsilon=u\)，故 \(u\) 局部光滑。对光滑函数在中心作二阶 Taylor 展开，圆周平均为 \(u(a)+r^2\Delta u(a)/4+o(r^2)\)。假设使二阶系数为零，故 \(\Delta u=0\)。
\end{proof}
径向磨光同时说明平均值性质为何能恢复所有阶导数：不是从连续性直接求导，而是先将函数识别成它的光滑卷积。

\subsection{最大值原理}
\label{b3:subsec:150103}

\begin{theorem}[强最大值原理]\label{b3:thm:15-maximum}
区域 \(D\) 上的实值调和函数若在内部取得全域最大值或最小值，则为常数。若 \(D\) 有界、\(u\in C(\overline D)\) 且在 \(D\) 调和，则最大值和最小值均可在边界取得。
\end{theorem}
\begin{proof}
设最大值为 \(M=u(a)\)。每个以 \(a\) 为中心的内部圆周上都有 \(u\le M\)，而平均等于 \(M\)。连续非负函数 \(M-u\) 的积分为零，故它在圆周恒零。所有充分小的半径合起来说明 \(u=M\) 的集合开；它由连续性又是闭集。连通性使其等于 \(D\)。最小值情形应用于 \(-u\)。有界域闭包紧，最大值存在；若极值只在内部出现，则上面的结论使函数处处相等，边界亦取得该值。
\end{proof}
有界性与边界条件各有作用。上半平面上的 \(u(x+\mathrm{i}y)=y\) 在实轴为零，却不恒零；没有无穷远控制，边界最大值原理和 Dirichlet 解的唯一性都会失败。

\subsection{Harnack 不等式}
\label{b3:subsec:150104}

为了比较不同点的正调和函数，先建立圆盘表示。在 \(|z|<R\) 记
\[
P_R(z,t)=\frac{R^2-|z|^2}{|Re^{\mathrm{i}t}-z|^2}.
\]
它是 \((Re^{\mathrm{i}t}+z)/(Re^{\mathrm{i}t}-z)\) 的实部，故关于 \(z\) 调和；几何级数表明其圆周平均为一。若 \(\varphi\) 连续，积分 \((2\pi)^{-1}\int P_R(z,t)\varphi(Re^{\mathrm{i}t})\,\mathrm{d}t\) 在内部调和。靠近边界一点时，远离该点的核一致趋零，近处由 \(\varphi\) 的一致连续性控制，所以积分连续取边值 \(\varphi\)。与 \(u\) 比较并用最大值原理，便得到 \(u\) 的 Poisson 表示。

\begin{theorem}[Harnack]\label{b3:thm:15-harnack}
若 \(u\ge0\) 在 \(B(0,R)\) 调和，则 \(|z|=r<R\) 时
\[
\frac{R-r}{R+r}\,u(0)\le u(z)\le\frac{R+r}{R-r}\,u(0).
\]
\end{theorem}
\begin{proof}
先在 \(r<S<R\) 的圆盘用上述表示。核满足 \((S-r)/(S+r)\le P_S(z,t)\le(S+r)/(S-r)\)，对非负边值积分后利用平均值性质得到不等式，再令 \(S\uparrow R\)。
\end{proof}
有限条相交圆盘构成的链把估计传到区域中任意固定紧集。这给出 Harnack 收敛原理：递增调和函数列若在一点有界，就局部一致收敛到调和函数；否则处处趋于正无穷。证明时把差函数代入链式估计，中心处的 Cauchy 性即变成紧集上的一致 Cauchy 性，再用平均值性质的逆命题。


平均值公式还给出内部导数估计。取 \(\overline{B(a,2r)}\subset D\)，在 \(B(a,r)\) 上用半径 \(2r\) 的 Poisson 表示。对核关于 \(z\) 的各阶导数求界，分母距离至少 \(r\)，得到
\[
\sup_{B(a,r)}|\partial^\alpha u|\le C_\alpha r^{-|\alpha|}\sup_{B(a,2r)}|u|.
\]
常数只依赖多重指标，因为把变量缩放到单位圆盘之后，核导数在半径二分之一的闭圆盘一致有界。估计同时控制导数与局部尺度；区域靠近边界时可用的 \(r\) 变小，不能把同一个内部常数当作边界导数界。

\ProseParagraphBreak

若调和函数族在每个紧集上统一有界，以上估计给导数的统一界和等度连续性。对一列这样的函数，先在第一个紧集抽出一致收敛子列，再逐层抽取并取对角列；所得局部一致极限保留平均等式，因此调和。对更高导数重用同一估计，可得到局部各阶导数收敛。这里的紧性机制与全纯正规族相似，但使用的是实值核和平均，而不是复曲线积分。

\ProseParagraphBreak

具体计算可以检验这些结论。边值 \(\varphi(t)=1+\cos t\) 的单位圆盘延拓为 \(1+r\cos\theta\)，非负，且在边界点 \(-1\) 为零。内部仍严格正，因为正调和函数只有在处处零时才可在内部为零。相反，\(u=x\) 在直径上有很多内部零点，却不违背强最小值原理，因为它在圆盘中不非负。不能把“调和函数无内部最小值”误读成“调和函数不能有内部零点”。

\ProseParagraphBreak

Harnack 原理中一点有界的量词也值得区分。若递增列 \(u_n\) 全部调和但不非负，可以先减 \(u_1\)，使差列非负；在一点上界便由 Harnack 链传到任何固定紧集。若没有单调性，一点有界不足以控制全列，例如 \(u_n(z)=n\operatorname{Re}z\) 在原点恒零，在别处发散。真正传播的是非负差函数的中心值，而不是任意函数的一个点值。


\ProseParagraphBreak

有界的整个调和函数也满足 Liouville 性质。设 \(|u|\le M\) 于平面，内部一阶导数估计在任意点、任意半径给 \(|\nabla u(a)|\le CM/R\)，令 \(R\to\infty\) 得梯度零，所以常数。实际上只要求单侧有界即可：若 \(u\ge0\)，对固定 \(z\) 在任意大圆盘应用 Harnack，不等式两端的系数均趋一，故 \(u(z)=u(0)\)。若只要求在某条直线有界，结论不成立，\(u(x+\mathrm{i}y)=y\) 已给反例。全平面的范围在扩大圆盘这一证明中被明确使用。

\ProseParagraphBreak

圆盘上的正调和函数还有边界测度表示，不要求其连续达到边界。对 \(0<s<R\)，把非负测度 \(u(se^{\mathrm{i}t})\mathrm{d}t/(2\pi)\) 放在参数圆周，它的总质量恒为 \(u(0)\)。紧圆周上的有限测度族具有弱收敛子列；此处可用连续函数可数稠密集上的对角抽取构造正泛函，再用正泛函表示得到极限测度 \(\nu\)。固定 \(|z|<R\)，在 \(s>|z|\) 的圆盘使用 Poisson 表示，并让子列 \(s\uparrow R\)，核关于参数一致收敛，得到
\[
u(z)=\int_0^{2\pi}P_R(z,t)\,\mathrm{d}\nu(t),\qquad \nu([0,2\pi))=u(0).
\]
傅里叶展开确定了 \(\nu\) 的全部三角矩，再由三角多项式稠密性确定测度，故表示唯一。这是边界数据由函数扩大到测度的一种自然方式。

\ProseParagraphBreak

特别地，取 \(\nu\) 为一点质量，就得到 \(P_R(z,t_0)\)。它在几乎所有边界点趋零，却在 \(Re^{\mathrm{i}t_0}\) 趋于无穷，内部仍严格正。因而“边界几乎处处为零”不等于经典连续边值为零；若不规定解的函数类，不能从几乎处处边值推出唯一性。Poisson 核的这个例子会在 Hardy 空间中再次出现，届时可积边界代表和测度奇异部分必须分别处理。


\ProseParagraphBreak

正调和函数的导数还可用自身点值控制。若 \(u\ge0\) 调和于 \(B(a,2r)\)，Harnack 给 \(\sup_{B(a,r)}u\le C u(a)\)，再在较小圆盘使用一阶内部估计，得到 \(|\nabla u(a)|\le C u(a)/r\)。当 \(u(a)>0\)，除以自身值可得 \(|\nabla\log u(a)|\le C/r\)。所以正调和函数在一个内部尺度上不能变化得任意快；常数只依赖尺度比而不依赖函数大小。若允许变号，\(u_n=x/n+y\) 等函数在某个点可有零值而梯度非零，便不可能由自身值给同样的界。

\ProseParagraphBreak

平均值与面积测度的归一化在缩放下保持一致。将 \(z=a+r\zeta\) 后，面积变为 \(r^2\mathrm{d}A(\zeta)\)，Laplace 算子变为 \(r^{-2}\Delta_\zeta\)，圆周长度变为 \(r\mathrm{d}s\)。因此圆周平均用 \(2\pi r\) 除，面积平均用 \(\pi r^2\) 除；分布基本解中的点质量则在对应变量替换时保持单位总质量。这些尺度检查会在后面的 Green 核、Riesz 测度与零点计数中反复出现。只看公式的函数部分而忽略测度变量，是丢失二维常数与半径幂的常见原因。

% !TeX root = ../../Book3.tex
\section{调和共轭}
\label{b3:sec:1502}

在二维，寻找调和共轭等价于对一个闭微分形式积分。方程解决局部存在性，沿闭曲线的积分则决定局部对象能否拼成一个全局函数。

\subsection{Cauchy–Riemann 方程}
\label{b3:subsec:150201}

若希望 \(u+\mathrm{i}v\) 全纯，必须满足 \(v_x=-u_y\)、\(v_y=u_x\)。所需的一形式为 \(\omega=-u_y\,\mathrm{d}x+u_x\,\mathrm{d}y\)，其外微分为 \(\Delta u\,\mathrm{d}x\wedge\mathrm{d}y\)。所以调和方程恰好消去局部积分的路径依赖。这里的符号由 Cauchy--Riemann 方程确定，改成 \(u_y\,\mathrm{d}x-u_x\,\mathrm{d}y\) 会得到共轭的负号。

\subsection{局部调和共轭}
\label{b3:subsec:150202}

\begin{proposition}\label{b3:prop:15-local-conjugate}
实值调和函数在每点的某个圆盘邻域具有调和共轭，且共轭只差实常数。
\end{proposition}
\begin{proof}
以圆盘中心 \((a,b)\) 为起点，令
\[
v(x,y)=-\int_a^x u_y(t,b)\,\mathrm{d}t+\int_b^y u_x(x,t)\,\mathrm{d}t.
\]
得到 \(v_y=u_x\)，又由 \(u_{xx}=-u_{yy}\) 得 \(v_x=-u_y(x,b)-\int_b^yu_{yy}(x,t)\,\mathrm{d}t=-u_y(x,y)\)。因此 \(u+\mathrm{i}v\) 全纯，\(v\) 也调和。两个共轭的偏导相同，差在连通圆盘上为常数。
\end{proof}
例如 \(u=x^2-y^2\) 对应 \(v=2xy\)；\(u=\log|z|\) 的局部共轭为某个辐角分支。后一例在绕原点一周后增加 \(2\pi\)，已显示全局障碍。

\subsection{单连通区域}
\label{b3:subsec:150203}

\begin{theorem}\label{b3:thm:15-global-conjugate}
区域 \(D\) 上的实值调和函数 \(u\) 有单值调和共轭，当且仅当每条分段光滑闭曲线 \(\gamma\subset D\) 均满足 \(\int_\gamma\omega=0\)。特别地，单连通区域上一定存在共轭。
\end{theorem}
\begin{proof}
若 \(\omega=\mathrm{d}v\)，闭曲线积分为零。反之，从固定点到 \(z\) 积分 \(\omega\)；两条路径之差是闭曲线，故积分独立于路径，局部微分得到上一小节的偏导。单连通时，每条闭曲线可缩到一点；把同伦像用有限个局部原函数圆盘覆盖、细分参数方格，相邻边积分相消，故积分在同伦中不变。
\end{proof}
去心圆盘上的 \(\log|z|\) 给出 \(\omega=\mathrm{d}\arg z\)，单位圆上积分为 \(2\pi\)，所以单连通条件不能随意删去。

\subsection{全纯函数}
\label{b3:subsec:150204}

如果只需要全纯函数的模，可以先构造局部共轭。不同共轭的差为常数，故局部函数 \(e^{u+\mathrm{i}v}\) 在交集只差模一常数；它们不一定拼成单值函数。若所有周期落在 \(2\pi\mathbb Z\)，指数却能消除差异，得到全局无零点全纯函数。

\begin{proposition}\label{b3:prop:15-log-modulus}
区域上无零点全纯函数 \(f\) 的 \(\log|f|\) 调和。反之，调和函数 \(u\) 能写成 \(\log|f|\)，其中 \(f\) 无零点全纯，当且仅当 \(-u_y\,\mathrm{d}x+u_x\,\mathrm{d}y\) 的所有周期属于 \(2\pi\mathbb Z\)。
\end{proposition}
\begin{proof}
局部对数使第一结论成立，并使必要性化为辐角沿闭曲线的总变化。充分性由上面的路径积分定义 \(e^{u+\mathrm{i}\int\omega}\)，不同路径改变整数倍 \(2\pi\)，故指数相同；局部表示保证全纯。
\end{proof}


可直接从实部计算全纯函数的例子能帮助核对共轭符号。设
\(u=x^3-3xy^2+xy\)。先算 \(\Delta u=0\)，再由 \(v_y=u_x=3x^2-3y^2+y\) 积分得
\(v=3x^2y-y^3+y^2/2+h(x)\)。代入 \(v_x=-u_y=6xy-x\)，得到 \(h'(x)=-x\)，故
\[
v=3x^2y-y^3+\frac{y^2-x^2}{2}+C,
\quad f(z)=z^3-\frac{\mathrm{i}}2z^2+\mathrm{i}C.
\]
这个计算同时检查了调和条件、共轭的相容性和纯虚常数的唯一自由度。若一开始给出的 \(u\) 不调和，例如 \(x^2+y^2\)，两偏导方程会在混合导数处矛盾；这比形式地积分其中一条方程更能发现问题。

\ProseParagraphBreak

不同边界孔的周期可以各自贡献。在去掉 \(a_1,\ldots,a_k\) 的区域上，函数
\(u(z)=\sum c_j\log|z-a_j|\) 调和。沿只绕 \(a_j\) 一周的曲线，周期为 \(2\pi c_j\)。因此全局单值共轭存在当且仅当这些基本周期全部零；全局无零点全纯函数的对数模存在则要求 \(c_j\in\mathbb Z\)。后一情况由 \(f=\prod(z-a_j)^{c_j}\) 明确实现，在当前去点区域内负指数不会引入内部极点。两种存在问题的差异来自指数函数的周期，不是来自不同的调和方程。

\ProseParagraphBreak

单连通性是保证所有调和函数都有共轭的条件，却不是某个特定函数有共轭的必要条件。环域上 \(u=\operatorname{Re}z\) 有共轭 \(\operatorname{Im}z\)，而 \(\log|z|\) 没有。对一个给定 \(u\)，应检查它的闭形式周期；不能只看到域有孔就拒绝任何共轭构造。


\ProseParagraphBreak

局部共轭还使调和函数具有实解析性。在小圆盘写 \(u=\operatorname{Re}f\)，全纯函数的 Taylor 级数局部一致收敛，因此
\[
u(a+re^{\mathrm{i}\theta})=c_0+\sum_{n\ge1}r^n(A_n\cos n\theta+B_n\sin n\theta).
\]
这里 \(c_0\) 为实数，系数由 \(f\) 的复 Taylor 系数拆成实部、虚部得到。每个非恒定项本身调和，且圆周平均零。若 \(u\) 在某个非空开集恒零，则其局部全纯函数实部恒零，Cauchy--Riemann 方程使它为纯虚常数；沿重叠圆盘传播后，\(u\) 在连通区域恒零。这个唯一延拓结论比最大值原理要求的单侧极值条件不同：前者使用开集上的全部零值，后者使用一点极值。

\ProseParagraphBreak

若两个调和函数在一条曲线上相等，唯一延拓却未必适用。\(u=x\) 与零函数在虚轴相等但在两侧不同。若额外给出该直线上的法向导数相等，在具有解析延拓的局部情形，Taylor 系数才能被同时限制。只由曲线上的函数值不能恢复二维的全部局部数据；全纯函数沿有聚点的零集恒等定理不能直接挪给实值调和函数。

\ProseParagraphBreak

Harnack 不等式还把局部正性传遍连通区域。固定紧集 \(K\Subset D\) 和一点 \(z_0\in D\)，用有限条圆盘链把 \(z_0\) 连到覆盖 \(K\) 的小圆盘；每一步比较的常数只依赖圆盘的半径比例。相乘后得到 \(\sup_Ku\le C_{K,z_0}u(z_0)\)，对所有正调和函数使用同一个常数。连通性在这里转化为有限次局部比较，而非直接给一个全域常数。若区域无界或紧集逐渐碰到边界，常数一般随之增大；圆盘 Poisson 核靠近其边界质量所在点时的增长就是一个具体检验。

% !TeX root = ../../Book3.tex
\section{次调和函数}
\label{b3:sec:1503}

全纯函数在零点处的对数模趋于负无穷，不能作为通常的调和函数处理。把平均等式放宽为不等式，并允许这个无穷值，便能保留最大值原理和零点信息。

\subsection{次调和函数}
\label{b3:subsec:150301}

\begin{definition}\label{b3:def:15-subharmonic}
设 \(D\) 为区域，\(u:D\to[-\infty,\infty)\) 上半连续、局部可积且不恒为 \(-\infty\)。若每个内部闭圆盘满足
\[
u(a)\le\frac1{2\pi}\int_0^{2\pi}u(a+r e^{\mathrm{i}t})\,\mathrm{d}t,
\]
则称 \(u\) 为\term[cidiaohe-hanshu]{次调和函数}[subharmonic function]。圆周积分取扩展意义。函数 \(-u\) 称为\term[chaodiaohe-hanshu]{超调和函数}[superharmonic function]。
\end{definition}
\(C^2\) 函数次调和当且仅当 \(\Delta u\ge0\)：必要性用二阶平均展开，充分性用圆周平均导数的面积公式。因此 \(|z|^2\) 次调和而非调和，\(-|z|^2\) 超调和。局部可积条件保证后面的分布运算有定义，也排除处处无穷的退化情形。

\subsection{平均值不等式}
\label{b3:subsec:150302}

圆盘上的比较描述同样有用。若 \(h\) 在圆盘调和、连续到边界，且 \(u\le h\) 于边界，则 \(u\le h\) 于内部。证明可先对 \(u-h+\varepsilon|z|^2\) 使用严格平均不等式：它不能有内部最大值，随后令 \(\varepsilon\downarrow0\)。一般上半连续边值用连续上逼近处理。反过来，连续边值的 Poisson 解与 \(u\) 比较，就给中心平均不等式。

\begin{proposition}\label{b3:prop:15-subharmonic-operations}
有限个次调和函数的最大值次调和；非负线性组合只要不恒为负无穷也次调和。若 \(u_j\downarrow u\) 且极限局部可积，则 \(u\) 次调和。
\end{proposition}
\begin{proof}
取最大值的函数在中心选出达到最大的一项，将其平均用所有项的最大值控制。和的情形由积分线性性。下降极限上半连续；在固定圆周上减去局部共同上界，再用单调收敛，得到平均不等式。局部可积保证其为允许的对象。
\end{proof}
任意上确界一般不保持上半连续；这正是 Perron 法必须处理而不能略去的一点。

\subsection{最大值原理}
\label{b3:subsec:150303}

\begin{proposition}[次调和函数的最大值原理]\label{b3:prop:15-subharmonic-maximum}
次调和函数若在区域内部取得有限最大值，则恒等于该值。若 \(D\) 有界且每个边界点满足 \(\limsup_{z\to\xi}u(z)\le0\)，则 \(u\le0\) 于 \(D\)。
\end{proposition}
\begin{proof}
在最大点，圆周平均与中心值相同。若圆周某处小于最大值，上半连续性给一段严格下降的弧，使平均严格下降，矛盾。故最大值集合既开又闭。边界条件加有限覆盖使某个边界邻域中 \(u\le\varepsilon\)；若内部仍有值大于 \(\varepsilon\)，最大值将在余下的紧集取得，第一结论迫使常数，和边界矛盾。令 \(\varepsilon\downarrow0\)。
\end{proof}

\subsection{\texorpdfstring{\(\log|f|\)}{log|f|}}
\label{b3:subsec:150304}

\begin{proposition}\label{b3:prop:15-log-subharmonic}
非恒零全纯函数 \(f\) 的 \(\log|f|\) 次调和，零点处取 \(-\infty\)。
\end{proposition}
\begin{proof}
局部写 \(f(z)=(z-a)^m g(z)\)，\(g\) 无零点。函数 \(\log|g|\) 调和，\(\log|z-a|\) 局部可积。后者以连续函数 \(\tfrac12\log(|z-a|^2+\varepsilon^2)\) 下降逼近；其 Laplace 算子为 \(2\varepsilon^2/(|z-a|^2+\varepsilon^2)^2\ge0\)。于是它次调和；加上调和函数保持不等式。局部性质通过圆盘比较描述给全域性质。
\end{proof}
这一证明保留零点重数，后面会把 \(m\) 解释为 Laplace 测度在 \(a\) 的质量。不能在零点任意补一个有限值；那会破坏平均不等式；仅仅保持上半连续性并不足够。


严格平均不等式的比较证明可以避免对 \(-\infty\) 点求导。设圆盘 \(B\) 上 \(u\le h\) 于边界，若内部存在 \(u-h>0\)，选充分小 \(\varepsilon>0\)，使 \(u-h+\varepsilon(|z|^2-R^2)\) 仍在某个内部点为正，而边界非正。上半连续性使它在闭圆盘取得有限正最大值。以最大点为中心取一个内部小圆，\(|z|^2\) 的圆周平均比中心多 \(r^2\)，所以中心值严格小于平均，和最大性矛盾。该证明只使用上半连续性、平均不等式及一个严格次调和修正，适用于对数奇性。

\ProseParagraphBreak

在圆盘边界上，上半连续且有上界的数据可由连续函数从上逼近。具体可先将数据截断到 \([-N,M]\)，再用小邻域上确界和平滑分片构造连续上界；截断后的单调逼近回到扩展值。Poisson 核非负且总质量一，使这种上逼近与积分相容。这个过程说明比较描述和圆周平均描述为什么一致，也解释了为什么不能用任意的点态修改代替上半连续代表。

\ProseParagraphBreak

一个常见构造是凸函数作用于调和函数。若 \(u\) 调和、\(\Phi:\mathbb R\to\mathbb R\) 凸且连续，则 Jensen 不等式与平均值性质给
\[
\Phi(u(a))\le\frac1{2\pi}\int_0^{2\pi}\Phi(u(a+re^{\mathrm{i}t}))\,\mathrm{d}t.
\]
故 \(\Phi\circ u\) 次调和。取 \(\Phi(t)=|t|\)、\(e^t\) 可得到许多不是调和的函数；在 \(C^2\) 情况，\(\Delta\Phi(u)=\Phi''(u)|\nabla u|^2\) 显示非调和性集中在凸性的二阶项。凸性不能删，例如 \(\Phi(t)=-t^2\) 与 \(u=x\) 给 Laplace 为负。

\ProseParagraphBreak

两个实部组合成复全纯函数时，\(|f|^p\) 对每个 \(p>0\) 次调和。在无零点局部，用调和函数 \(\log|f|\) 与凸函数 \(e^{pt}\) 即得；在零点处，函数连续、中心值零，平均非负，故平均不等式直接成立。这个结论不需要 \(p\ge1\)，因为这里使用的是指数作用于对数模，而不是直接声称 \(t\mapsto|t|^p\) 在所有实数上凸。

% !TeX root = ../../Book3.tex
\section{分布观点}
\label{b3:sec:1504}

调和与次调和的差别还可以用一个正测度表达。它记录函数偏离平均等式的来源，而不要求函数具有通常的二阶导数。

\subsection{分布意义下的 Laplace 算子}
\label{b3:subsec:150401}

对于 \(u\in L^1_{\mathrm{loc}}(D)\)，定义 \(\langle\Delta u,\varphi\rangle=\int_Du\Delta\varphi\,\mathrm{d}A\)，其中 \(\varphi\in C_c^\infty(D)\)。分布配对不加复共轭，二阶积分分部使符号为正。
\begin{proposition}[对数基本解]\label{b3:prop:15-log-fundamental}
在复平面分布意义下，\(\Delta\log|z|=2\pi\delta_0\)。
\end{proposition}
\begin{proof}
删去半径 \(\varepsilon\) 的圆盘，在余域用 Green 公式。小圆的余域外法向朝向原点，边界项为
\[
\int_0^{2\pi}\bigl(-\varepsilon\log\varepsilon\,\partial_r\varphi(\varepsilon e^{\mathrm{i}t})+\varphi(\varepsilon e^{\mathrm{i}t})\bigr)\,\mathrm{d}t.
\]
第一项趋零，第二项趋 \(2\pi\varphi(0)\)。被删圆盘上 \(\log|z|\) 的可积性保证原积分极限成立。
\end{proof}

\subsection{正分布}
\label{b3:subsec:150402}

实分布 \(T\) 对每个非负测试函数取非负值时，称为\term[zheng-fenbu]{正分布}[positive distribution]。取 \(\chi\in C_c^\infty(D)\)，\(\chi=1\) 于紧集 \(K\) 邻域，则支集在 \(K\) 的实测试函数满足
\[
|T(\varphi)|\le T(\chi)\cdot\|\varphi\|_\infty.
\]
这是由 \(-\|\varphi\|_\infty\chi\le\varphi\le\|\varphi\|_\infty\chi\) 得到的。光滑逼近因此把 \(T\) 唯一延拓到连续紧支集函数，并保持正性。第一卷定理11.4.1（正泛函的 Riesz--Markov 表示）给出局部有限正 Radon 测度。这里调用的是表示定理，不把分布的高阶连续性直接误当作测度性。

\subsection{Riesz 测度}
\label{b3:subsec:150403}

\begin{proposition}\label{b3:prop:15-riesz-measure}
次调和函数 \(u\) 的分布 \(\Delta u\) 为正分布。故有唯一局部有限正测度
\(\mu_u=(2\pi)^{-1}\Delta u\)，称为其\term[riesz-cedu]{Riesz 测度}[Riesz measure]。
\end{proposition}
\begin{proof}
径向非负磨光 \(u_\varepsilon\) 在缩小的区域中仍满足平均不等式；它光滑，故 \(\Delta u_\varepsilon\ge0\)。又 \(u_\varepsilon\to u\) 于局部 \(L^1\)，对非负测试函数取极限得到 \(\int u\Delta\varphi\ge0\)。上一小节把该正分布识别为测度，唯一性由测试函数稠密性得到。
\end{proof}
例如 \(u=|z|^2\) 的测度为 \((2/\pi)\mathrm{d}A\)，而 \(u=\log|z-a|\) 的测度为 \(\delta_a\)。连续分布与点质量都可以出现，不能把 Riesz 测度只理解为离散零点。

\subsection{势论解释}
\label{b3:subsec:150404}

对光滑紧支集非负密度 \(\rho\)，对数势
\[
V(z)=\int_{\mathbb C}\log|z-w|\cdot\rho(w)\,\mathrm{d}A(w)
\]
局部可积，且 \(\Delta V=2\pi\rho\)。换序由紧支集上的对数核局部可积保证。因而若 \(\mu_u=\rho\mathrm{d}A\)，\(u-V\) 的分布 Laplace 算子为零。

分布调和函数的光滑性可以这样恢复：其磨光 \(h_\varepsilon\) 调和；在一个内部圆盘上再用固定径向核 \(\eta\)，平均值性质给 \(h_\varepsilon=h_\varepsilon*\eta\)。局部 \(L^1\) 极限使右端连同每个阶导数一致收敛，得到光滑调和代表 \(h=h*\eta\)。这是 \term[weyl-yinli]{Weyl 引理}[Weyl's lemma]在平面上的证明。

势决定非调和部分，调和函数决定不被测度看见的自由度。在一般 Radon 密度情形，对数势可能在薄集上为负无穷，必须另行选择次调和代表；这里不把局部可积函数的任意零测修改说成同一个点态次调和函数。


正分布的零阶估计还说明其局部有限性。取 \(K\Subset D\)，所有 \(0\le\varphi\le1\)、支集在 \(K\) 的测试函数均满足 \(T(\varphi)\le T(\chi)\)。表示测度在 \(K\) 的质量因此不超过同一有限常数。不同紧集选择不同 \(\chi\)；不能从局部有限性推出全域测度有限。

\ProseParagraphBreak

例如 \(u(z)=|z|^2\) 的 Riesz 测度在每个有界紧集有限，却在全平面质量无穷。\(u=\log|z|\) 的测度总质量一，但是函数在无穷远不有界。测度有限性与函数有界性是不同性质，后面 Green 势的边界控制不能只根据其中一项猜出。

\ProseParagraphBreak

正分布一定是测度，任意分布却不一定如此。点质量的导数 \(\partial_x\delta_0\) 对测试函数的作用是 \(-\partial_x\varphi(0)\)。可取支集缩小、最大模不超过一而原点导数任意大的光滑函数，所以它不满足紧集上的零阶最大范数界，不能由局部有限测度表示。它也不正：非负测试函数在原点的一阶导数可有不同符号，只要原点不是极值点。正性在测度识别中承担的正是把所有阶连续性降为零阶控制的作用。

\ProseParagraphBreak

对于无零点全纯函数，\(\mu_{\log|f|}=0\)，但函数的增长可以很大。例如 \(e^{z^3}\) 的对数模为 \(\operatorname{Re}z^3\)，调和且测度零。因此 Laplace 测度只记录势的源，并不记录全部调和自由度。要确定一个函数，还须给边界数据、增长条件或某种归一化；这也是下一章边值问题与再下一章 Green 最小性各自出现的原因。


\ProseParagraphBreak

对一般有限正测度也可以局部构造对数势。设 \(\mu\) 支集于紧集 \(K\)，定义 \(V(z)=\int_K\log|z-w|\,\mathrm{d}\mu(w)\)，允许负无穷。对任意有界集 \(L\)，
\[
\sup_{w\in K}\int_L|\log|z-w||\,\mathrm{d}A(z)<\infty,
\]
因为单位距离以内的负部由 \(2\pi\int_0^1r|\log r|\,\mathrm{d}r\) 控制，外面的正部由 \(K,L\) 的有界性控制。Tonelli 因而给 \(V\in L^1_{\mathrm{loc}}\)，尤其它不能在整个开集恒为负无穷。分布换序及对数基本解得到 \(\Delta V=2\pi\mu\)。

\ProseParagraphBreak

这个积分本身提供了适当的点态代表。对局部 \(z\) 变量，核有共同有限上界；减去该上界后用 Fatou 引理，得到 \(\limsup V(z_j)\le V(z)\)，故上半连续。每个核的圆周平均不小于中心值，对正测度积分后仍保持该不等式；若负部的换序有疑问，可先以光滑对数核逼近，再用下降极限处理。因此 \(V\) 次调和，而不只是一个 Laplace 正的分布等价类。

\ProseParagraphBreak

给定次调和 \(u\)，在两个同心内部圆盘 \(B'\Subset B\Subset D\) 上取 \(\mu=\mu_u|_{\overline B}\)，上述势满足 \(\Delta(u-V)=0\) 于 \(B\)，Weyl 引理给调和函数 \(h\)，于是 \(u=V+h\) 几乎处处。要恢复处处等式，可利用次调和函数的面积平均在半径趋零时回到自身：平均不等式给下界，上半连续性给上界，所以平均极限等于点值，包括值为负无穷的情形。双方几乎处处相等使所有面积平均相等，最终得到 \(u=V+h\) 处处成立。这是局部 Riesz 分解；把全部测度放入势时，远处边界与增长还须另外控制，本证明只承诺内部局部范围。


\begin{chaptersummary}
平均值性质把调和方程转为局部比较。Poisson 核的正性给出 Harnack 控制，核的可微性给出内部导数估计；这些估计又使局部有界的调和函数族具有紧致性。使用正调和函数的表示或 Harnack 链时，正性承担着控制全部函数值的作用，不能仅以某一点的有界性代替。

\ProseParagraphBreak
调和共轭在小圆盘上总能构造，全局存在却要检验闭曲线周期。周期全为零时，共轭可取成单值函数；周期只要求落在 \(2\pi\mathbb Z\) 时，指数已经单值，因而得到全局的全纯对数模。环域中的对数项展示了这两个条件的差别。

\ProseParagraphBreak
次调和函数以平均不等式保留对数奇性，其 Riesz 测度是 \((2\pi)^{-1}\Delta u\)。正分布的零阶估计保证这一对象确实是局部有限测度，分布判据则从几乎处处等价类中选出唯一次调和代表。对数势具有指定的 Riesz 测度，Weyl 引理把剩余的零 Laplace 分布恢复为调和函数，从而得到局部分解。测度决定非调和部分，边界条件、周期或增长限制还要控制调和部分；后续边值问题将进一步讨论这些条件如何配合。
\end{chaptersummary}

\BookExercises

% 八题均为自拟或对原自拟题的实质改编；新增训练分别涉及非径向周期、径向凸性、单侧增长、曲线源、圆环势、边界唯一性、高阶极值和势的收敛。解答只使用本章及此前工具，未附会第三方题源。
\Needspace{6\baselineskip}
\begin{exercise}[由圆周均值识别环域中的障碍]\label{b3:ex:15-draft-01}
设 \(0<\rho<R\)，实函数 \(u\) 在环域 \(A=\{\rho<|z|<R\}\) 调和，记
\[
 m(r)\coloneqq\frac1{2\pi}\int_0^{2\pi}u(re^{\mathrm it})\,\mathrm dt.
\]
证明对任意 \(\rho<r_1<r_2<R\)，商
\[
 c\coloneqq\frac{m(r_2)-m(r_1)}{\log(r_2/r_1)}
\]
与所选半径无关，并且存在 \(A\) 上的单值全纯函数 \(F\)，使 \(u=c\log|z|+\operatorname{Re}F\)。据此判别何时存在无零点全纯函数 \(f\) 满足 \(\log|f|=u\)，写出所有这样的 \(f\)；给定整数 \(n\ge2\)，再判别这样的 \(f\) 何时具有 \(A\) 上的全纯 \(n\) 次根。
\end{exercise}
\begin{solution}
在内部圆周上对平均求导，并用极坐标中的调和方程，得到
\[
 (rm'(r))'
 =\frac1{2\pi}\int_0^{2\pi}(ru_{rr}+u_r)\,\mathrm dt
 =-\frac1{2\pi r}\int_0^{2\pi}u_{tt}\,\mathrm dt=0.
\]
所以 \(m(r)=c\log r+d\)，所给商确实与半径无关。沿正向圆周积分共轭形式 \(\omega=-u_y\,\mathrm dx+u_x\,\mathrm dy\)，有
\[
 \int_{|z|=r}\omega
 =\int_0^{2\pi}r u_r(re^{\mathrm it})\,\mathrm dt=2\pi c.
\]
环域的闭曲线周期由这一圆周周期生成。函数 \(u-c\log|z|\) 的共轭形式在生成圆周上的积分为零，因而全部周期为零。由全局共轭判据，它是某个单值全纯函数 \(F\) 的实部。这也证明分解中的实系数 \(c\) 唯一；\(F\) 只差纯虚常数。

若 \(u=\log|f|\)，对数模的周期条件要求 \(2\pi c\in2\pi\mathbb Z\)，即 \(c\in\mathbb Z\)。反过来，当 \(c\) 为整数时，\(z^c e^F\) 在环域上全纯且无零点，对数模正是 \(u\)。任何另一解与它的商模恒为1，故为模一常数。因此全部解为
\[
 f(z)=e^{\mathrm i\beta}z^c e^{F(z)},\quad \beta\in\mathbb R.
\]
若 \(f=g^n\)，则 \(g\) 也无零点，且 \(\log|g|=u/n\)。将已经证明的判据用于 \(u/n\)，得到必要条件 \(c/n\in\mathbb Z\)。此条件也充分，因为
\[
 g(z)=e^{\mathrm i\beta/n}z^{c/n}e^{F(z)/n}
\]
给出一个全纯根，其余根只差一个 \(n\) 次单位根。判别只取决于圆周均值的对数斜率，并不要求原函数径向对称。
\end{solution}

\Needspace{6\baselineskip}
\begin{exercise}[径向次调和性与凸函数的斜率]\label{b3:ex:15-draft-02}
设 \(I=(\log\rho,\log R)\)，其中 \(0<\rho<R\)，且 \(\psi:I\to\mathbb R\) 连续。在环域 \(A=\{\rho<|z|<R\}\) 上定义 \(u(z)=\psi(\log|z|)\)。证明 \(u\) 次调和当且仅当 \(\psi\) 凸。凸时令 \(\nu=\mathrm d\psi'_+\) 为其右导数确定的 Stieltjes 测度，证明对 \(\varphi\in C_c(A)\)，
\[
 \int_A\varphi\,\mathrm d\mu_u
 =\frac1{2\pi}\int_I\int_0^{2\pi}
       \varphi(e^{t+\mathrm i\theta})\,\mathrm d\theta\,\mathrm d\nu(t).
\]
求 \(\mu_u(\{|z|=r\})\) 及 \(\mu_u(\{r<|z|<s\})\)，其中 \(\rho<r<s<R\)，并说明导数跳跃如何产生圆周上的质量。
\end{exercise}
\begin{solution}
对实测试函数 \(\varphi\in C_c^\infty(A)\)，记
\[
 \Phi(t,\theta)=\varphi(e^{t+\mathrm i\theta}),\quad
 \overline\Phi(t)=\frac1{2\pi}\int_0^{2\pi}\Phi(t,\theta)\,\mathrm d\theta.
\]
令 \(r=e^t\) 作极坐标换元。面积因子与 Laplace 算子的径向因子相消，角向二阶导数积分为零，于是
\[
 \langle\Delta u,\varphi\rangle
 =\int_I\int_0^{2\pi}\psi(t)
       (\Phi_{tt}+\Phi_{\theta\theta})\,\mathrm d\theta\,\mathrm dt
 =2\pi\int_I\psi(t)\overline\Phi''(t)\,\mathrm dt.
\]
若 \(u\) 次调和，取 \(\varphi(z)=\chi(\log|z|)\)，其中 \(0\le\chi\in C_c^\infty(I)\)，便得 \(\int_I\psi\chi''\ge0\)。这表示一维分布 \(\psi''\) 非负。在任意内部区间上磨光 \(\psi\)，所得光滑函数的二阶导数非负，所以为凸函数；连续性使磨光局部一致收敛到 \(\psi\)，凸不等式随极限保留，故 \(\psi\) 凸。

反过来，有限凸函数局部绝对连续，右导数 \(\psi'_+\) 单调且局部有界，其 Stieltjes 测度 \(\nu\) 为正测度。对光滑紧支集 \(\chi\)，两次分部积分分别按绝对连续函数和 Stieltjes 积分理解，给出
\[
 \int_I\psi\chi''\,\mathrm dt
 =-\int_I\psi'_+\chi'\,\mathrm dt
 =\int_I\chi\,\mathrm d\nu.
\]
代回前式，得到 \(\Delta u\ge0\) 及题中的测度公式，先对光滑测试函数成立，再由测度唯一性延拓到连续紧支集函数。由于 \(u\) 连续，局部磨光回到其原有点值；分布判据因此说明这个 \(u\) 本身次调和。

Stieltjes 测度在一点的质量是右、左极限之差，故
\[
 \mu_u(\{|z|=r\})
 =\psi'_+(\log r)-\psi'_-(\log r),
\]
而开环域的质量为
\[
 \mu_u(\{r<|z|<s\})
 =\psi'_-(\log s)-\psi'_+(\log r).
\]
前一质量按角度均匀分布在圆周上。例如斜率在 \(t_0\) 向上跳 \(a>0\)，就产生总质量 \(a\) 的均匀圆周测度，圆周半径为 \(e^{t_0}\)。若 \(\psi\in C^2\)，则相应的面积密度为 \(\psi''(\log|z|)/(2\pi|z|^2)\)。这两个描述允许光滑密度与圆周质量同时出现。
\end{solution}

\Needspace{6\baselineskip}
\begin{exercise}[单侧增长也限制调和函数的次数]\label{b3:ex:15-draft-03}
设实调和函数 \(u\) 定义在整个复平面，且存在 \(C>0\)、实数 \(\alpha\ge0\)，使
\[
 u(z)\le C(1+|z|)^\alpha.
\]
证明 \(u\) 是次数至多 \(\lfloor\alpha\rfloor\) 的调和多项式；这里没有假设 \(-u\) 满足同样的上界。进一步，若整数 \(N\ge1\) 满足
\[
 \max_{|z|=R}u(z)=o(R^N)\quad(R\to\infty),
\]
证明次数至多 \(N-1\)，并给出说明两种次数界都不能改进的例子。
\end{exercise}
\begin{solution}
先把单侧上界转为较小圆盘上的双侧上界。设 \(M_R\ge0\) 且 \(u\le M_R\) 于 \(B(0,2R)\)。正调和函数 \(M_R-u\) 的 Harnack 不等式给出，对 \(|z|\le R\)，
\[
 M_R-u(z)\le3(M_R-u(0)).
\]
结合 \(u(z)\le M_R\)，得到
\[
 \sup_{|z|\le R}|u(z)|\le2M_R+3|u(0)|.
\]
第一种假设下可取 \(M_R=C(1+2R)^\alpha\)。固定 \(a\in\mathbb C\)，令整数 \(k=\lfloor\alpha\rfloor+1\)。当 \(R\) 充分大时，\(a\) 附近半径与 \(R\) 成固定比例的圆盘包含在 \(B(0,R)\) 内，内部导数估计因此给出，对每个 \(|\beta|=k\)，
\[
 |\partial^\beta u(a)|
 \le C_\beta R^{-k}\bigl(2C(1+2R)^\alpha+3|u(0)|\bigr).
\]
因 \(k>\alpha\)，令 \(R\to\infty\) 得所有 \(k\) 阶导数为零。Taylor 公式于是表明 \(u\) 是次数至多 \(k-1\) 的多项式；它原来满足的调和方程仍成立。

第二种假设下，取
\[
 M_R=\max\left(0,\max_{|z|=2R}u(z)\right)=o(R^N).
\]
最大值原理使它也控制整个 \(B(0,2R)\)。同一 Harnack 估计给 \(\sup_{|z|\le R}|u(z)|=o(R^N)\)，其中常数项因 \(N\ge1\) 可并入小 \(o\)。这次对所有 \(N\) 阶导数使用内部估计，极限为零，故次数至多 \(N-1\)。

对第一种界，\(u(z)=\operatorname{Re}z^{\lfloor\alpha\rfloor}\) 达到所许次数；\(\lfloor\alpha\rfloor=0\) 时取非零常数。对第二种界，\(\operatorname{Re}z^{N-1}\) 满足小 \(o\) 条件且次数恰为 \(N-1\)。单侧条件能够控制次数，关键在于每个大圆盘上都可以先减去一个共同上界，再使用正性。
\end{solution}

\Needspace{6\baselineskip}
\begin{exercise}[双曲线上的 Riesz 测度]\label{b3:ex:15-draft-04}
在复平面上设
\[
 h(z)=\operatorname{Re}(z^2)-1,\quad u(z)=\max(h(z),0),\quad
 \Gamma=\{x^2-y^2=1\}.
\]
直接用测试函数证明
\[
 \mathrm d\mu_u=\frac{|z|}{\pi}\,\mathrm ds_\Gamma,
\]
其中 \(\mathrm ds_\Gamma\) 为 \(\Gamma\) 上的弧长测度。进一步求 \(\mu_u(B(0,R))\) 对所有 \(R>0\) 的显式表达式，并说明 \(\Gamma\) 的两个顶点处有无点质量。
\end{exercise}
\begin{solution}
调和函数 \(h\) 与零函数的最大值次调和。记 \(E=\{h>0\}\)；在其边界 \(\Gamma\) 上，\(h=0\) 且 \(\nabla h=(2x,-2y)\ne0\)。对 \(\varphi\in C_c^\infty(\mathbb C)\)，在 \(E\) 与足够大的圆盘之交上使用 Green 公式。外圆上的测试函数及其导数均为零，故
\[
 \langle\Delta u,\varphi\rangle
 =\int_E h\Delta\varphi\,\mathrm dA
 =\int_\Gamma(h\partial_n\varphi-\varphi\partial_n h)\,\mathrm ds.
\]
\(E\) 的外法向朝向 \(h\) 减小的一侧，所以 \(\partial_n h=-|\nabla h|=-2|z|\)。由 \(h=0\) 得
\[
 \langle\Delta u,\varphi\rangle
 =\int_\Gamma2|z|\varphi\,\mathrm ds,
\]
除以 \(2\pi\) 即得所求测度。

当 \(R\le1\)，开圆盘不与 \(\Gamma\) 相交，质量为零。当 \(R>1\)，两支曲线可分别参数化为
\[
 x=\pm\sqrt{1+y^2},\quad |y|<Y,\quad
 Y=\sqrt{\frac{R^2-1}{2}}.
\]
沿任一支有
\[
 |z|\,\mathrm ds
 =\frac{1+2y^2}{\sqrt{1+y^2}}\,\mathrm dy
 =\mathrm d\bigl(y\sqrt{1+y^2}\bigr).
\]
把两支的积分相加，得到
\[
 \mu_u(B(0,R))=
 \begin{cases}
 0,&0<R\le1,\\[2pt]
 \dfrac2\pi\sqrt{R^4-1},&R>1.
 \end{cases}
\]
两个顶点 \((\pm1,0)\) 都是光滑曲线上的普通点，密度在其附近连续；弧长测度对单点为零，所以没有额外点质量。这里的非调和部分集中在整个曲线界面，不能只从函数不光滑的几个可见位置猜测点质量。
\end{solution}

\Needspace{6\baselineskip}
\begin{exercise}[均匀圆环的对数势]\label{b3:ex:15-draft-05}
设 \(0<a<b\)，令
\[
 \mathrm d\mu(w)=\frac{\mathbf1_{\{a<|w|<b\}}}{\pi(b^2-a^2)}\,\mathrm dA(w),\quad
 V(z)=\int_{\mathbb C}\log|z-w|\,\mathrm d\mu(w).
\]
求 \(V\) 在 \(|z|\le a\)、\(a<|z|<b\)、\(|z|\ge b\) 三个区域的显式表达式。证明势在中空圆盘内为常数，且在两条边界圆周处一阶导数连续；据此直接核对 \(\mu_V=\mu\)，包括边界圆周上是否出现质量。
\end{exercise}
\begin{solution}
记 \(r=|z|\)、\(D=b^2-a^2\)。旋转不变性使 \(V\) 只依赖 \(r\)。当 \(r\ne s\) 时，提出较大的半径，再对无零点函数的全纯对数实部使用平均值性质，可得
\[
 \frac1{2\pi}\int_0^{2\pi}\log|r-se^{\mathrm it}|\,\mathrm dt
 =\log\max(r,s).
\]
例如 \(r<s\) 时，余项为 \(\log|1-(r/s)e^{-\mathrm it}|\)，平均为零。对数核的局部可积性允许极坐标换序；参数 \(s=r\) 这单个半径不影响面积积分。因此
\[
 V(r)=\frac2D\int_a^b s\log\max(r,s)\,\mathrm ds.
\]
利用 \(\int s\log s\,\mathrm ds=(s^2/2)\log s-s^2/4\)，分别计算得到
\[
 V(r)=
 \begin{cases}
 \dfrac{b^2\log b-a^2\log a}{D}-\dfrac12,&0\le r\le a,\\[6pt]
 \dfrac{b^2\log b-a^2\log r-(b^2-r^2)/2}{D},&a<r<b,\\[6pt]
 \log r,&r\ge b.
 \end{cases}
\]
第一段显示中空圆盘内的势为常数，第三段则等于总质量为1的原点点质量在外部产生的势；在中间的有质量区域，二者都不能代替实际的势。

三段在端点取值一致，径向导数分别为
\[
 V'(r)=
 \begin{cases}
 0,&0<r<a,\\[2pt]
 \dfrac{r^2-a^2}{Dr},&a<r<b,\\[6pt]
 \dfrac1r,&r>b.
 \end{cases}
\]
在 \(a\) 处中间导数趋于0，在 \(b\) 处趋于 \(1/b\)，故 \(V\) 在整个平面为 \(C^1\) 函数。在开圆环内，\(V''+V'/r=2/D\)，而在其余两区为零。分区使用 Green 公式时，函数值和法向一阶导数的边界项都相消，所以分布 Laplace 算子没有附加的圆周项。因此
\[
 \frac1{2\pi}\Delta V
 =\frac{\mathbf1_{\{a<|z|<b\}}}{\pi D}\,\mathrm dA=\mu.
\]
\end{solution}

\Needspace{6\baselineskip}
\begin{exercise}[几乎处处边值何时能够确定调和函数]\label{b3:ex:15-draft-06}
设实函数 \(u\) 在 \(\mathbb D\) 调和，且对几乎每个 \(t\in[0,2\pi)\)，有 \(u(re^{\mathrm it})\to0\) 当 \(r\uparrow1\)。若存在 \(1<p<\infty\)，使
\[
 \sup_{0<r<1}\frac1{2\pi}\int_0^{2\pi}|u(re^{\mathrm it})|^p\,\mathrm dt<\infty,
\]
证明 \(u\equiv0\)。再给出满足同样径向边值条件、所有圆周 \(L^1\) 范数均等于1的非零正调和函数，说明把 \(p>1\) 改成 \(p=1\) 后结论不成立。
\end{exercise}
\begin{solution}
以 \(\mathrm dm=\mathrm dt/(2\pi)\) 表示归一化圆周测度，令 \(\|u(re^{\mathrm it})\|_{L^p(m)}\le M\)。固定 \(\varepsilon>0\)，设
\[
 E_r=\{t:|u(re^{\mathrm it})|>\varepsilon\}.
\]
几乎处处收敛使 \(m(E_r)\to0\)。在 \(E_r\) 上用 Hölder 不等式，在其补集上用 \(\varepsilon\) 控制，得到
\[
 \int|u(re^{\mathrm it})|\,\mathrm dm
 \le\varepsilon+M\,m(E_r)^{1-1/p}.
\]
先令 \(r\uparrow1\)，再令 \(\varepsilon\downarrow0\)，证明圆周 \(L^1\) 范数趋于零。正的指数 \(1-1/p\) 在这一步不可省略。

固定 \(z\in\mathbb D\)。对 \(|z|<r<1\) 使用半径 \(r\) 的 Poisson 表示：
\[
 u(z)=\int_0^{2\pi}
       \frac{r^2-|z|^2}{|re^{\mathrm it}-z|^2}
       u(re^{\mathrm it})\,\frac{\mathrm dt}{2\pi}.
\]
当 \(r\) 充分接近1时，核由 \((r+|z|)/(r-|z|)\) 一致控制。右端的绝对值因圆周 \(L^1\) 范数趋零而趋零，故 \(u(z)=0\)。任意性给 \(u\equiv0\)。

反例为边界点1的 Poisson 核
\[
 v(z)=\operatorname{Re}\frac{1+z}{1-z}
     =\frac{1-|z|^2}{|1-z|^2}.
\]
它正调和，\(v(0)=1\)，且对 \(t\not\equiv0\pmod{2\pi}\)，有 \(v(re^{\mathrm it})\to0\)。由正性和均值性质，\(\int|v(re^{\mathrm it})|\,\mathrm dm=1\) 对每个 \(r\) 成立。圆周质量越来越集中在角度0附近，正是这种集中使有界的 \(L^1\) 范数无法推出 \(L^1\) 收敛。
\end{solution}

\Needspace{6\baselineskip}
\begin{exercise}[正调和函数的高阶方向导数]\label{b3:ex:15-draft-07}
设 \(u\ge0\) 在 \(\mathbb D\) 调和，\(u(0)=1\)，整数 \(n\ge2\)。对 \(\theta\in\mathbb R\)，定义方向导数
\[
 \mathrm D_\theta^n u(0)
 \coloneqq\left.\frac{\mathrm d^n}{\mathrm ds^n}u(se^{\mathrm i\theta})\right|_{s=0}.
\]
证明 \(|\mathrm D_\theta^n u(0)|\le2n!\)，并用边界表示测度的支集分别刻画正、负两种取等情形。若再有 \(u(z)=1+O(|z|^n)\) 当 \(z\to0\)，证明正向取等 \(\mathrm D_\theta^n u(0)=2n!\) 唯一确定
\[
 u(z)=\operatorname{Re}\frac{1+e^{-\mathrm in\theta}z^n}{1-e^{-\mathrm in\theta}z^n}.
\]
\end{exercise}
\begin{solution}
令 \(\nu\) 为 \(u\) 的边界表示测度，它是单位圆周上的概率测度。Poisson 核在每个内部紧圆盘上一致展开为
\[
 \frac{1-|z|^2}{|\zeta-z|^2}
 =1+2\operatorname{Re}\sum_{k=1}^\infty z^k\overline\zeta^{\,k},
 \quad |\zeta|=1.
\]
积分并沿给定方向求 \(n\) 阶导数，得到
\[
 \mathrm D_\theta^n u(0)
 =2n!\int_{\partial\mathbb D}
       \operatorname{Re}(e^{\mathrm in\theta}\overline\zeta^{\,n})\,\mathrm d\nu(\zeta).
\]
被积函数介于 \(-1\) 与1之间，所求界立即成立。

积分等于1，当且仅当 \(\nu\) 集中在被积函数等于1的地方；这是因为非负连续函数 \(1-\operatorname{Re}(e^{\mathrm in\theta}\overline\zeta^{\,n})\) 的积分必须为零。因此正向取等当且仅当
\[
 \operatorname{supp}\nu\subseteq
 \{\zeta\in\partial\mathbb D:\zeta^n=e^{\mathrm in\theta}\}.
\]
同理，负向取等当且仅当支集包含在 \(\{\zeta:\zeta^n=-e^{\mathrm in\theta}\}\) 中。两种集合都各有 \(n\) 个点，其上的任意概率权重都给出取等函数，所以常数确实达到，且一般并不只有一个取等函数。

在正向取等情形，写 \(\omega=e^{2\pi\mathrm i/n}\)，则
\[
 \nu=\sum_{j=0}^{n-1}a_j\delta_{e^{\mathrm i\theta}\omega^j},
 \quad a_j\ge0,\quad\sum_j a_j=1.
\]
条件 \(u(z)-1=O(|z|^n)\) 使前 \(n-1\) 个复矩 \(\int\overline\zeta^{\,k}\,\mathrm d\nu\) 全为零：若其中首个非零矩的阶数为 \(k<n\)，Poisson 展开在某个方向就有非零的 \(k\) 阶首项，与该增长条件矛盾。于是
\[
 \sum_{j=0}^{n-1}a_j\omega^{-jk}=0\quad(1\le k<n).
\]
连同 \(k=0\) 时和为1，对这些等式作有限 Fourier 反演，或直接用单位根的正交关系，可得 \(a_j=1/n\)。均匀权重使全部非 \(n\) 的倍数阶矩消失，而第 \(\ell n\) 阶矩为 \(e^{-\mathrm i\ell n\theta}\)。因此
\[
 u(z)=1+2\operatorname{Re}\sum_{\ell=1}^\infty
                 (e^{-\mathrm in\theta}z^n)^\ell
     =\operatorname{Re}\frac{1+e^{-\mathrm in\theta}z^n}{1-e^{-\mathrm in\theta}z^n}.
\]
它确实正调和并满足全部要求，唯一性亦由权重的确定得到。
\end{solution}

\Needspace{6\baselineskip}
\begin{exercise}[测度弱收敛与对数势的两种收敛]\label{b3:ex:15-draft-08}
设 \(\mu_j,\mu\) 为支集都包含在同一紧集 \(K\subset\mathbb C\) 中的有限正测度，且对每个 \(f\in C(K)\)，有 \(\int f\,\mathrm d\mu_j\to\int f\,\mathrm d\mu\)。记其对数势为
\[
 V_j(z)=\int_K\log|z-w|\,\mathrm d\mu_j(w),\quad
 V(z)=\int_K\log|z-w|\,\mathrm d\mu(w).
\]
证明 \(V_j\to V\) 于 \(L^1_{\mathrm{loc}}(\mathbb C)\)。再构造一列概率测度，使上述弱收敛成立，但在一个固定点 \(z_0\) 处，\(V_j(z_0)\) 均有限且不趋于 \(V(z_0)\)。
\end{exercise}
\begin{solution}
对常函数1应用弱收敛，得到总质量有共同有限上界 \(C\)，并可把 \(\mu(K)\) 也包括在这个上界内。固定紧圆盘 \(L\)，对 \(M>0\) 定义下截断核
\[
 k_M(z,w)=\max(\log|z-w|,-M),
\]
在 \(z=w\) 时取值 \(-M\)。它在 \(L\times K\) 上连续有界。记相应的截断势为 \(V_{j,M}\)、\(V_M\)；对每个固定 \(z\)，弱收敛给 \(V_{j,M}(z)\to V_M(z)\)，共同的质量上界又提供关于 \(z\in L\) 的一致常数界。支配收敛定理因此给
\[
 \|V_{j,M}-V_M\|_{L^1(L)}\longrightarrow0
 \quad\text{对每个固定的 }M.
\]

截断误差可以独立于测度的集中位置来控制。Tonelli 定理和整个小圆盘上的估计给出
\begin{align*}
 \|V_{j,M}-V_j\|_{L^1(L)}
 &\le C\int_{|s|<e^{-M}}(-M-\log|s|)\,\mathrm dA(s)\\
 &=2\pi C\int_0^{e^{-M}}r(-M-\log r)\,\mathrm dr
 =\frac{\pi C}{2}e^{-2M}.
\end{align*}
同一估计适用于 \(V_M-V\)。这些势可能在个别点为 \(-\infty\)，但局部可积性保证此处的 \(L^1\) 计算有意义。由三角不等式，先令 \(j\to\infty\)，再令 \(M\to\infty\)，得到 \(\|V_j-V\|_{L^1(L)}\to0\)。紧圆盘任意，故有局部 \(L^1\) 收敛。

点态结论的反例可保持总质量始终为1。对 \(j\ge2\)，令
\[
 \mu_j=\left(1-\frac1j\right)\delta_1
              +\frac1j\delta_{e^{-j}},\quad \mu=\delta_1.
\]
所有支集均包含在实线段 \([0,1]\) 中；任意连续测试函数有界，而第二项的质量趋零，所以 \(\mu_j\) 弱收敛到 \(\mu\)。在固定观察点 \(z_0=0\)，却有
\[
 V_j(0)=\left(1-\frac1j\right)\log1
                  +\frac1j\log(e^{-j})=-1,\quad V(0)=0.
\]
每个 \(V_j(0)\) 都有限，但不收敛到极限势的点值。越来越靠近观察点的少量质量可以被对数奇性放大；核的积分尾界仍足以保证局部 \(L^1\) 收敛，却不控制这个固定点的值。
\end{solution}
